\subsection{Complements on the relative separable algebraic closure}

THEOREM 1 (Zariski). --- Let K be a field, L an extension of K, \(K_{1}\) the relative separable algebraic closure of K in L (V, p.~44), \(K_{2}\) the relative algebraic closure of K in L (V, p.~19) and \((\mathbf{X}_{i})_{i\in I}\) a family of indeterminates. Then \(\mathrm{K}_{1}(\mathbf{X}_{i})_{i\in I}\) is the relative separable algebraic closure of \(\mathrm{K}(\mathbf{X}_{i})_{i\in I}\) in \(\mathrm{L}(\mathbf{X}_{i})_{i\in I}\) and \(\mathrm{K}_{2}(\mathbf{X}_{i})_{i\in I}\) is the relative algebraic closure of \(\mathrm{K}(\mathbf{X}_{i})_{i\in I}\) in \(\mathrm{L}(\mathbf{X}_{i})_{i\in I}\) .

\begin{enumerate}
\def\labelenumi{\Alph{enumi})}
\tightlist
\item
  Suppose that E is a field and F an extension field of E and that every element of F which is separable algebraic over E belongs to E. Let u be an element of \(F(X)\) which is separable algebraic over \(E(X)\) ; we shall show that u belongs to \(E(X)\) . There exist in \(F(X)\) two relatively prime polynomials P and Q such that u = P/Q and we may take Q to be monic. Let S be the finite subset of F consisting of the coefficients of P and Q, \(F_{0} = E(S)\) and let \(\Delta\) be an E-derivation of \(\mathbf{F}_0\) into \(\mathbf{F}_0\) . Let \(D\) be the derivation of \(\mathrm{F}_0(\mathbf{X})\) into itself which agrees with \(\Delta\) on \(\mathbf{F}_0\) and maps \(X\) to 0 (V, p.~128, Prop. 3).
\end{enumerate}

Since \(u \in F_0(X)\) is separable algebraic over \(E(X)\) and \(D\) is zero on \(E(X)\) , we have \(D(u) = 0\) (V, p.~129, Prop. 4), whence \(D(P).Q = P.D(Q)\) . Since \(P\) and \(Q\) are relatively prime, we conclude that \(Q\) divides \(D(Q)\) (IV, p.~13, Cor. 4). Now \(Q\) may be written in the form

\[
\mathrm{Q} (\mathrm{X}) = \mathrm{X} ^ {n} + a _ {1} \mathrm{X} ^ {n - 1} + \dots + a _ {n - 1} \mathrm{X} + a _ {n}
\]

with \(a_1, \ldots, a_n\) in \(\mathbf{F}_0\) ; since \(\mathbf{D}(x) = 0\) , we thus have

\[
\mathrm{D} (\mathrm{Q}) = \Delta (a _ {1}) \mathrm{X} ^ {n - 1} + \dots + \Delta (a _ {n - 1}) \mathrm{X} + \Delta (a _ {n})\tag{1}
\]

whence \(\deg D(Q) < \deg Q\) . Since \(Q\) divides \(D(Q)\) , this is possible only if \(D(Q) = 0\) . But then \(D(P) = 0\) because \(D(P) \cdot Q = P \cdot D(Q)\) . Now (1) and a similar formula for \(P\) show that \(\Delta\) annihilates the set \(S\) of coefficients of \(P\) and \(Q\) , whence \(\Delta = 0\) , because \(F_0 = E(S)\) . By \(V\) , p.~135, Cor. 2, the finitely generated extension \(F_0\) of \(E\) is thus separable algebraic; by the hypotheses on \(E\) and \(F\) we have \(F_0 = E\) , so finally \(u \in E(X)\) .

\begin{enumerate}
\def\labelenumi{\Alph{enumi})}
\setcounter{enumi}{1}
\item
  Suppose now that E is relatively algebraically closed in the extension field F and denote by p the characteristic exponent of E. Let u be an element of \(F(X)\) which is algebraic over \(E(X)\) . There exists an integer \(f \geqslant 0\) such that \(v = u^{p^f}\) is separable algebraic over \(E(X)\) (V, p.~44, Prop. 13). By A) we thus have \(v \in E(X)\) . There exists a unique representation of u in the form P/Q with relatively prime polynomials P and Q in F{[}X{]} and Q monic; we have a similar decomposition \(v = P_1 / Q_1\) with \(P_1\) and \(Q_1\) relatively prime in E{[}X{]} and \(Q_1\) monic. It follows that \(P_1 / Q_1 = P^{p^f} / Q^{p^f}\) ; the polynomials \(P^{p^f}\) and \(Q^{p^f}\) are relatively prime in F{[}X{]} (IV, p.~13, Cor. 6), like \(P_1\) and \(Q_1\) , and \(Q^{p^f}\) is monic. It follows that \(P^{p^f} = P_1 \in E[X]\) and \(Q^{p^f} = Q_1 \in E[X]\) . Therefore the coefficients of P and Q are p-radical over E, and hence belong to E because E is relatively algebraically closed in F. We thus have \(P \in E[X]\) , \(Q \in E[X]\) and finally \(u \in E(X)\) . This proves \(E(X)\) to be relatively algebraically closed in F(X).
\item
  We use the notation of Theorem 1. Since \(\mathbf{K}_1\) is a separable algebraic extension of \(\mathbf{K}\) , the extension \(\mathbf{K}_1(\mathbf{X}_i)_{i\in I}\) of \(\mathbf{K}(\mathbf{X}_i)_{i\in I}\) is algebraic and separable (V, p.~39, Prop. 6). Moreover, every element of \(L\) which is algebraic and separable over \(\mathbf{K}_1\) belongs to \(\mathbf{K}_1\) (V, p.~44, Prop. 13, a)). Let \(J\) be a finite subset of \(I\) ; by an immediate induction on the cardinal of \(J\) we deduce from \(A\) ) that every element of \(\mathrm{L}(\mathbf{X}_i)_{i\in J}\) which is algebraic and separable over \(\mathbf{K}(\mathbf{X}_i)_{i\in J}\) belongs to \(\mathbf{K}_1(\mathbf{X}_i)_{i\in J}\) . Finally let \(u\) be an element of \(\mathrm{L}(\mathbf{X}_i)_{i\in I}\) algebraic and separable over \(\mathbf{K}(\mathbf{X}_i)_{i\in I}\) ; there exists a finite subset \(J\) of \(I\) such that \(u\) belongs to \(\mathrm{L}(\mathbf{X}_i)_{i\in J}\) and is algebraic and separable over \(\mathbf{K}(\mathbf{X}_i)_{i\in J}\) ; by what has been said, \(u\) belongs to \(\mathbf{K}_1(\mathbf{X}_i)_{i\in J}\) and a fortiori to \(\mathbf{K}_1(\mathbf{X}_i)_{i\in I}\) .
\end{enumerate}

We have thus deduced from \(A\) that \(\mathrm{K}_1(\mathbf{X}_i)_{i\in I}\) is the relative separable closure of \(\mathrm{K}(\mathbf{X}_i)_{i\in I}\) in \(\mathrm{L}(\mathbf{X}_i)_{i\in I}\) ; in the same way it follows from \(B\) that \(\mathrm{K}_2(\mathbf{X}_i)_{i\in I}\) is the relative algebraic closure of \(\mathrm{K}(\mathbf{X}_i)_{i\in I}\) in \(\mathrm{L}(\mathbf{X}_i)_{i\in I}\) .

